\chapter{The prototype platform}\label{ch:platform}

This chapter describes the memory and devices of the prototype platform, as implemented identically by
the emulator and the RTL simulation harness (\texttt{common/platform.h}). Unlike the instruction set,
the platform is discovered at run time from the platform table, so software should not hard-code it.

\section{System memory map}\label{sec:memmap}

System memory (T3, space 3) is 64~MiB. The architecture gives it no structure: at reset core 0 holds
one capability for all of it. The platform fixes one region, the framebuffer window, and the boot ROM
fixes another, its load area (Section~\ref{sec:boot}).

\begin{table}[htbp]
\centering\small
\caption{System memory (T3) of the prototype platform}\label{tab:memmap}
\begin{tabularx}{\linewidth}{@{}l l X@{}}
\toprule
\textbf{Range} & \textbf{Size} & \textbf{Use}\\\midrule
\hex{0000000}--\hex{0EFFFFF} & 15 MiB & general memory; handed to the OS in \reg{c0}\\
\hex{0F00000}--\hex{0FFFFFF} & 1 MiB & boot load area: the OS code, data and far images, immutable after boot\\
\hex{1000000}--\hex{2FFFFFF} & 32 MiB & general memory; handed to the OS in \reg{c13}\\
\hex{3000000}--\hex{3FFFFFF} & 16 MiB & framebuffer window: the display scans out its first
  $\textit{height} \times \textit{stride}$ bytes (Section~\ref{sec:fbctl})\\
\bottomrule
\end{tabularx}
\end{table}

The framebuffer window is ordinary memory in every architectural respect: it is reached by transfers
and atomics through capabilities cut from the memory root, it can be split, given, shared and placed
revocable, and the part of it the display does not scan out is free for any use. The
cluster store (T2, space 2) is 256~KiB with no fixed regions.

\section{Device space}

Device space (space 4) is \hex{5000} bytes long. Core 0's device root capability \reg{c3} covers all of
it; software cuts narrower capabilities from copies of \reg{c3} with \insn{csplit}. Every port,
register and window starts at a multiple of 8, so a capability can be cut to exactly one of them
(Section~\ref{sec:capalign}).

\begin{table}[htbp]
\centering\small
\caption{Device space map}\label{tab:devmap}
\begin{tabularx}{\linewidth}{@{}l l l X@{}}
\toprule
\textbf{Offset} & \textbf{Length} & \textbf{Device} & \textbf{Ports, registers and windows}\\\midrule
\hex{0000} & \hex{100} & console & \hex{0000} output stream; \hex{0008} input stream; \hex{0010} break
                                     port; \hex{0018} break release\\
\hex{1000} & 320 & platform table & read-only window (\texttt{STPD})\\
\hex{2000} & \hex{300} & block device 0 & \hex{2000} identify window (64 B); \hex{2100} submission port;
                                          \hex{2200} completion port\\
\hex{3000} & \hex{20} & timer & \hex{3000} NOW; \hex{3008} FREQ; \hex{3010} DEADLINE; \hex{3018} WAIT\\
\hex{4000} & \hex{10} & framebuffer controller & \hex{4000} PRESENT; \hex{4008} VSYNC\\
\hex{4100} & \hex{10} & pointer & \hex{4100} EVENT; \hex{4108} STATE\\
\hex{4200} & \hex{20} & network interface & \hex{4200} TX; \hex{4208} RX; \hex{4210} STATUS;
                                            \hex{4218} CANCEL\\
\hex{4300} & \hex{18} & sound (optional) & \hex{4300} PCM; \hex{4308} RATE; \hex{4310} STATUS\\
\bottomrule
\end{tabularx}
\end{table}

Device transfers follow the rules of Section~\ref{sec:devxfer}: one byte per beat to consecutive
offsets; registers are 8 bytes, little-endian, latched when byte 0 is read and committed when byte 7 is
written; a transfer that starts at a stream port belongs to that stream. Reads from unassigned device
offsets return zero; writes to them are ignored. Six ports are \term{event ports}, whose reads stay
pending until the device has something to deliver (Section~\ref{sec:pendread}): console input, console
break, the timer's WAIT, the controller's VSYNC, the pointer's EVENT and the network interface's RX.

\section{Console}\label{sec:console}

\begin{description}[leftmargin=1em,style=nextline]
\item[Output (port \hex{0000})] Every byte of a \insn{tput} that starts at this port is output in order.
  Every core starts with \reg{c2}, a copyable write-only capability for exactly this port.
\item[Input (port \hex{0008})] Each byte of a \insn{tget} that starts at this port is the next input
  byte. The read is a pending device read: it waits until its bytes are available, so a \insn{tget} of
  $n$ bytes completes only after $n$ bytes have arrived, and meanwhile nothing else is held up. Pending
  reads are served oldest first. At end of input the port returns \hex{04} (EOT) for every further
  byte. Programs normally read one byte at a time.
\item[Break (port \hex{0010})] A \insn{tget} from this port completes when the user presses Ctrl-C,
  with the byte 3. Presses that arrive while no read is pending are counted and complete later reads at
  once, one press per read. Ctrl-C typed at the console is not delivered to the input port.
\item[Release (port \hex{0018})] Writing a byte to this port completes, with zeros, \emph{every} break
  read pending at that moment, on every core. A break read issued later waits for the next Ctrl-C or
  release. A write while no break read is pending has no effect.
\end{description}

The break port is how the console signals an event without an interrupt
(Section~\ref{sec:events}): a core that wants to be told about Ctrl-C keeps a one-byte read of the
break port in flight and waits for its tag, typically together with other event sources, with
\insn{tany}. A core retires its break read with a release before it halts (Section~\ref{sec:halt}).

On FPGA builds the console is a UART (8N1, 115200 baud). In the simulators it is the host's standard
output and standard input (or a file given with \texttt{-i}); when standard input is a terminal the
simulators pass keys through unechoed, the guest echoes them, and Ctrl-C goes to the break port. When
the simulators' graphics window has the keyboard focus, its keys feed the same input stream
(Section~\ref{sec:window}). The \texttt{-b} option of both simulators presses Ctrl-C at a given point
of the run for tests.

\section{Platform table}\label{sec:stpd}

The platform table describes the machine: memory tiers, cores and devices. It is a read-only window at
device offset \hex{1000}, 256 bytes long: a 32-byte header and seven 32-byte entries. All fields are
little-endian.

\begin{table}[htbp]
\centering\small
\caption{Platform table header (32 bytes)}\label{tab:stpdhdr}
\begin{tabular}{@{}r l l@{}}
\toprule
\textbf{Offset} & \textbf{Type} & \textbf{Field}\\\midrule
0 & char[4] & magic \texttt{"STPD"}\\
4 & u16 & version (1)\\
6 & u16 & entry size in bytes (32)\\
8 & u16 & number of entries (9)\\
10 & u16 & number of cores $N$\\
12 & u32 & total length of the table in bytes (320)\\
16 & u32 & checksum: sum of all bytes of the table, with this field taken as zero, mod $2^{32}$\\
20 & u16 & $\textit{CTAB}$: entries of every core's capability save table (64--256; 256 by default). 0 in an older table means 64\\
22 & u16 & $\textit{NXMAP}$: entries of every core's overlay map (4--16; 16 by default). 0 in an older table means 16\\
24 & u16 & \texttt{DLS\_SIZE} in KiB: every core's D-LS (a power of two, 4--64; 64 by default). 0 in an older table means 64\\
26 & u16 & \texttt{ILS\_SIZE} in KiB: every core's I-LS (the same range, independent of the D-LS). 0 in an older table means 64\\
28 & u8[4] & reserved, zero\\
\bottomrule
\end{tabular}
\end{table}

\begin{table}[htbp]
\centering\small
\caption{Platform table entry (32 bytes)}\label{tab:stpdent}
\begin{tabularx}{\linewidth}{@{}r l l X@{}}
\toprule
\textbf{Offset} & \textbf{Type} & \textbf{Field} & \textbf{Meaning}\\\midrule
0 & u16 & type & 1 memory, 2 console, 3 platform table, 4 block device, 5 framebuffer, 6 timer,
  7 pointer, 8 network interface, 9 sound\\
2 & u16 & index & instance number among entries of the same type\\
4 & u16 & space & far space of the resource (2, 3 or 4)\\
6 & u16 & node & \insn{cgive} node number, or \hex{ffff} if the device takes no capabilities\\
8 & u64 & info & type-specific (Table~\ref{tab:stpdinfo})\\
16 & u64 & base & base address in \textit{space}\\
24 & u64 & size & size in bytes\\
\bottomrule
\end{tabularx}
\end{table}

\begin{table}[htbp]
\centering\small
\caption{The \textit{info} word by entry type}\label{tab:stpdinfo}
\begin{tabularx}{\linewidth}{@{}r l X@{}}
\toprule
\textbf{Type} & \textbf{Name} & \textbf{\textit{info}}\\\midrule
1 & memory & 0\\
2 & console & bit 0 = has output, bit 1 = has input\\
3 & platform table & 1\\
4 & block device & capacity in 512-byte blocks\\
5 & framebuffer & bits 15:0 width in pixels; bits 31:16 height in pixels; bits 47:32 stride in bytes
  per row; bits 63:48 pixel format (1 = XRGB8888)\\
6 & timer & ticks per second (1000000)\\
7 & pointer & bits 7:0 buttons present (bit 0 left, 1 right, 2 middle); bits 15:8 wheels present
  (bit 8 vertical, bit 9 horizontal): \hex{0307}\\
8 & network interface & the MTU in bytes (1500)\\
9 & sound & samples the device's queue holds (16384)\\
\bottomrule
\end{tabularx}
\end{table}

The prototype publishes nine entries:

\begin{center}\small
\begin{tabular}{@{}l l l l l l l@{}}
\toprule
\textbf{\#} & \textbf{Type} & \textbf{Index} & \textbf{Space} & \textbf{Node} & \textbf{Base} & \textbf{Size}\\\midrule
0 & memory & 0 & 3 & --- & 0 & 64 MiB\\
1 & memory & 1 & 2 & --- & 0 & 256 KiB\\
2 & console & 0 & 4 & --- & \hex{0000} & \hex{100}\\
3 & platform & 0 & 4 & --- & \hex{1000} & 320\\
4 & block & 0 & 4 & $N$ & \hex{2000} & \hex{300}\\
5 & framebuffer & 0 & 3 & --- & \hex{3000000} & 16 MiB\\
6 & timer & 0 & 4 & --- & \hex{3000} & \hex{20}\\
7 & pointer & 0 & 4 & --- & \hex{4100} & \hex{10}\\
8 & network & 0 & 4 & --- & \hex{4200} & \hex{20}\\
\bottomrule
\end{tabular}
\end{center}

A machine with the sound device (Section~\ref{sec:sound}) publishes a tenth entry, of type 9, at
\hex{4300} with size \hex{18}; its table is 352 bytes long.

The framebuffer entry describes the framebuffer \emph{memory}: a region of system memory. The
framebuffer controller's two registers are at the fixed device offset \hex{4000} and have no entry of
their own.

\section{Timer}\label{sec:timer}

The timer is four 8-byte registers at device offset \hex{3000}. It counts \term{ticks}; one tick is one
microsecond.

\begin{table}[htbp]
\centering\small
\caption{Timer registers}\label{tab:timer}
\begin{tabularx}{\linewidth}{@{}l l l X@{}}
\toprule
\textbf{Offset} & \textbf{Name} & \textbf{Access} & \textbf{Function}\\\midrule
\hex{3000} & NOW & read & ticks since power-on\\
\hex{3008} & FREQ & read & ticks per second: 1000000\\
\hex{3010} & DEADLINE & read, write & an absolute tick count; 0 at reset\\
\hex{3018} & WAIT & read (event port) & the read completes when NOW $\ge$ DEADLINE and returns NOW as of
  its completion\\
\bottomrule
\end{tabularx}
\end{table}

\begin{itemize}
\item NOW is latched when its byte 0 is read, so an 8-byte read returns one consistent value.
\item A write to DEADLINE takes effect when its byte 7 is written. Pending WAIT reads are re-evaluated
  at once: a deadline in the past, in particular 0, completes them immediately.
\item There is one DEADLINE register. Several cores may have a WAIT pending at the same time; each
  completes when NOW reaches the single DEADLINE. A system that needs many timeouts gives the timer to
  one core, which keeps a queue of deadlines in software and always programs the earliest.
\item A WAIT read issued when NOW $\ge$ DEADLINE completes at once.
\item Writes to NOW, FREQ and WAIT are ignored.
\end{itemize}

Waiting for a point in time is therefore a transfer like any other event
(Section~\ref{sec:events}): write DEADLINE, start a read of WAIT on some tag, and wait for the tag with
\insn{tany}, together with whatever else the core is waiting for.

\begin{asm}
        la    t2, dlbuf                 ; the absolute deadline, 8 bytes in the D-LS
        li    t3, 0x10                  ; DEADLINE (offsets within the timer capability c5)
        li    t4, 8
        tput  tg2, t2, c5, t3, t4       ; takes effect when its byte 7 lands
        la    t2, wbuf
        li    t3, 0x18                  ; WAIT
        tget  tg1, t2, c5, t3, t4       ; pending until NOW >= DEADLINE; tg1 stays busy
        li    t0, (1 << 16) | (1 << 1)  ; a gift, or tag 1 idle
        tany  t1, t0                    ; sleep until either
\end{asm}

Reading the clock needs only \perm{R} on NOW and FREQ, so an operating system can hand every program
a read-only capability for exactly those 16 bytes and keep DEADLINE and WAIT to itself. StrataOS does
(Section~\ref{sec:os-time}).

\subsection{Time base}
\begin{description}[leftmargin=1em,style=nextline]
\item[Emulator] \texttt{-T real} (the default): NOW is host wall-clock microseconds since the start of
  the run. \texttt{-T virtual}: NOW is the number of machine rounds executed, divided by 100 (a nominal
  100~MHz machine in which every core executes one instruction per cycle); when every core is waiting
  for the timer, virtual time advances to the next deadline.
\item[RTL] A 64-bit counter of clock cycles divided by the parameter \texttt{CLK\_MHZ} (default 100).
  In simulation, \texttt{-T real} (the default) takes NOW from the host clock and \texttt{-T virtual}
  uses the cycle counter.
\end{description}
Runs whose output is compared between machines use \texttt{-T virtual}; even then the two machines'
clocks advance at different rates per instruction, so timer values are timing-visible
(Section~\ref{sec:determinism}).

\section{Framebuffer and display controller}\label{sec:fbctl}

The display has two parts: the \term{framebuffer}, which is memory, and the \term{framebuffer
controller}, which is two device registers.

\subsection{The framebuffer}
The framebuffer is the window \hex{3000000}--\hex{3FFFFFF} of system memory (16~MiB). The display scans
out its first $\textit{height} \times \textit{stride}$ bytes, row after row from the top.

\begin{description}[leftmargin=1em,style=nextline]
\item[Geometry] is a property of the machine, reported by the framebuffer entry of the platform table
  (Table~\ref{tab:stpdinfo}); software must not assume it. The default is 1280 $\times$ 800 pixels. Both
  simulators accept \texttt{-r W}x\texttt{H}, with a width that is a multiple of 8, from 320 $\times$ 200
  up to 2560 $\times$ 1600, which is the largest size the window holds.
\item[Pixels] are 4 bytes; the stride is $\textit{width} \times 4$ bytes. Format 1 is XRGB8888: a pixel is
  a little-endian 32-bit word with red in bits 23:16, green in bits 15:8 and blue in bits 7:0; in
  memory the bytes are B, G, R and one unused byte. A doubleword holds two pixels, the left one in its
  low half.
\item[Scanout is live.] The display continuously shows what is in the framebuffer memory. A write into
  it, by \insn{tput}, by \insn{tamo} or by a block-device read into a framebuffer capability, becomes
  visible at the next refresh with nothing else to do. Software that must not show a half-drawn frame
  draws elsewhere and copies.
\item[Ownership] is by capability, like all memory. Whoever holds a writable capability for a range of
  the framebuffer is the only one who can draw there; a renderer on $N$ cores splits the framebuffer
  capability into $N$ bands and gives one to each core (the \texttt{gfx} example,
  Section~\ref{sec:ex-gfx}), and a window system keeps the capability to itself and lends clients
  surfaces elsewhere (Section~\ref{sec:os-ws}).
\end{description}

\subsection{The controller}

\begin{table}[htbp]
\centering\small
\caption{Framebuffer controller registers}\label{tab:fbctl}
\begin{tabularx}{\linewidth}{@{}l l l X@{}}
\toprule
\textbf{Offset} & \textbf{Name} & \textbf{Access} & \textbf{Function}\\\midrule
\hex{4000} & PRESENT & write & ``this frame is complete'': the display is refreshed at once. The value
  written is ignored\\
\hex{4008} & VSYNC & read (event port) & the read completes at the next display refresh and returns the
  frame number\\
\bottomrule
\end{tabularx}
\end{table}

The display refreshes 60 times per second of timer time: a refresh occurs every 16667 ticks, at the
ticks that are multiples of 16667, and the frame number is NOW divided by 16667. A VSYNC read always
waits for the \emph{next} refresh, so a loop that reads VSYNC once per frame is paced at 60~Hz, and
several cores may wait at once.

PRESENT does not copy anything and is not needed for pixels to appear. It marks a point at which the
framebuffer holds a complete frame: the simulators repaint their window at once and record that frame
for the \texttt{-F} option (Section~\ref{sec:emu}), which is what tests compare. A program that
animates draws a frame, writes PRESENT and reads VSYNC.

\begin{asm}
        tput  tg0, s1, c9, zero, t1     ; PRESENT: 8 bytes at offset 0 of the controller capability
        li    t0, 8                     ; VSYNC
        tget  tg1, s1, c9, t0, t1       ; pending until the next refresh
        ld    a0, 0(s1)                 ; the frame number: this load waits for the refresh
\end{asm}

In the simulators the framebuffer appears in a window, opened the first time software writes into the
framebuffer or writes PRESENT (Section~\ref{sec:window}). The FPGA builds have no video output:
PRESENT has no effect there, and VSYNC still paces at 60~Hz of timer time.

\section{Pointer}\label{sec:pointer}

The pointer is a mouse or touchpad over the display, with three buttons and a wheel: two 8-byte
registers at device offset \hex{4100}.

\begin{table}[htbp]
\centering\small
\caption{Pointer registers}\label{tab:pointer}
\begin{tabularx}{\linewidth}{@{}l l l X@{}}
\toprule
\textbf{Offset} & \textbf{Name} & \textbf{Access} & \textbf{Function}\\\midrule
\hex{4100} & EVENT & read (event port) & the read completes with the next queued event\\
            &       & write & completes every EVENT read pending at that moment, on every core, with 0
  (the value written is ignored)\\
\hex{4108} & STATE & read & the pointer now (below); reading it changes nothing\\
\bottomrule
\end{tabularx}
\end{table}

An event is 8 bytes: bits 15:0 x and bits 31:16 y in framebuffer pixels; bits 39:32 the buttons
held (bit 0 left, 1 right, 2 middle); bits 47:40 the vertical wheel and bits 55:48 the horizontal
wheel, signed notches (positive: away from the user, and right); bits 63:56 the kind: 1 move, 2
button, 3 wheel, 4 enter, 5 leave. Every event carries the position and the buttons as they were. A
read released by a write returns 0, which is no kind. STATE holds x and y in the same places, the
buttons in bits 39:32, the number of events queued in bits 47:40, and in bit 56 whether the pointer is
over the display.

The device queues 16 events. Consecutive moves become one, carrying the last position; consecutive
wheel turns at the same place and with the same buttons add up; and when the queue is full the oldest
move makes room, so button and wheel events are never dropped for moves. The first move of a pointer
that was not over the display is an enter. A release lets software that is stopping (an operating
system halting the machine) complete the reads it left waiting, as the console's break release does.

The simulators feed the pointer from their window's mouse and hide the host cursor once software has
read the device, since software that uses it draws its own. For tests, \texttt{-P}~\textit{file} on both
simulators scripts it (Section~\ref{sec:emu}). The FPGA builds have no pointer: STATE reads 0 and an
EVENT read waits for a release.

\section{Network interface}\label{sec:netif}

The network interface carries raw IPv4 packets: there is no link-layer framing and there are no
hardware addresses below IP. It is four 8-byte ports at device offset \hex{4200}.

\begin{table}[htbp]
\centering\small
\caption{Network interface ports}\label{tab:netif}
\begin{tabularx}{\linewidth}{@{}l l l X@{}}
\toprule
\textbf{Offset} & \textbf{Name} & \textbf{Access} & \textbf{Function}\\\midrule
\hex{4200} & TX & write (stream) & one transfer is one packet: the bytes of a \insn{tput} that starts here
  are sent as one packet when the transfer ends\\
\hex{4208} & RX & read (event port) & completes when a packet is queued, with that packet\\
\hex{4210} & STATUS & read & bits 15:0 the length of the next queued packet (0: none); bits 23:16
  packets queued; bit 24 link up; bit 25 the far end is a user-mode NAT; bits 31:26 a host number the
  link suggests (0: none); bits 63:32 packets dropped\\
\hex{4218} & CANCEL & write & completes every RX read pending at that moment, on every core, with
  zeros (the value written is ignored)\\
\bottomrule
\end{tabularx}
\end{table}

A transfer to TX longer than the MTU (1500 bytes, the platform entry's \textit{info}) is dropped and
counted. An RX read takes the oldest queued packet: byte $i$ of the read is byte $i$ of the packet, the
bytes past the packet's end are zero, and the part of a packet past the read's length is lost; software
learns the length from the IPv4 header. The device queues 16 received packets, and one that finds the
queue full is dropped and counted. A CANCEL takes no packet; like the console's break release, it lets
software that is stopping complete the read it left waiting (StrataOS's kernel writes it when it halts
the machine).

What is at the other end of the link is chosen with the simulators' \texttt{-N} option
(Section~\ref{sec:emu}): nothing (the default: packets sent are dropped, nothing arrives, STATUS says
the link is down), a loopback plug (every packet sent is received back at once), a built-in test peer
that answers pings and echoes UDP datagrams to port~7, or a ``cable'' to another simulator made of host
UDP datagrams, one per packet, so that two simulated machines can talk TCP to each other. The FPGA
builds have no PHY yet: TX is dropped, RX waits for a CANCEL and STATUS reads 0.

\section{Sound}\label{sec:sound}

The sound device is a PCM output: a queue of mono, signed 16-bit samples that it plays at a set
rate. It is optional. A machine that has it lists it in the platform table (type 9); software that
finds no such entry goes without sound. The reference emulator has it when started with
\texttt{-a} or \texttt{-A}; the RTL does not, so there the three offsets are unassigned.

\begin{table}[htbp]
\centering\small
\caption{Sound registers}\label{tab:sound}
\begin{tabularx}{\linewidth}{@{}l l l X@{}}
\toprule
\textbf{Offset} & \textbf{Name} & \textbf{Access} & \textbf{Function}\\\midrule
\hex{4300} & PCM & write (stream port) & the transfer's bytes are samples, low byte first, appended
  to the queue; samples that find the queue full are dropped\\
\hex{4308} & RATE & read, write & samples per second, 8000 to 48000 (a value outside is clamped);
  11025 at power-on\\
\hex{4310} & STATUS & read & bits 31:0 samples queued and not yet played; bits 63:32 samples dropped
  since power-on\\
\bottomrule
\end{tabularx}
\end{table}

The queue holds 16384 samples and empties at RATE in timer time; an empty queue plays silence.
Software paces itself with STATUS: it keeps a fraction of a second queued and writes more as the
queue runs down. Samples written by several cores are queued in the order their transfers arrive.

\section{Block device}\label{sec:blkdev}

Block device 0 (the \texttt{STBD} protocol) stores 512-byte blocks. Its distinguishing feature is how it
is given access to memory: the device can touch far memory only through a capability that a core
\emph{gives} it. Handing over the buffer capability \emph{is} the DMA permission, and the device hands it
back when the operation is complete. While the device holds the buffer, no core can touch it, so DMA
protection holds by construction.

\subsection{Identify window}
\begin{center}\small
\begin{tabular}{@{}r l l@{}}
\toprule
\textbf{Offset} & \textbf{Type} & \textbf{Field}\\\midrule
0 & char[4] & magic \texttt{"STBD"}\\
4 & u16 & version (1)\\
6 & u16 & $\log_2$ block size (9)\\
8 & u64 & capacity in blocks\\
16 & u32 & maximum blocks per command (256)\\
20 & u32 & queue depth (4): submission and completion queues\\
24 & u8[40] & reserved, zero\\
\bottomrule
\end{tabular}
\end{center}

\subsection{Commands and completions}
A command is 32 bytes written with one \insn{tput} to the submission port; a completion is 16 bytes read
with one \insn{tget} from the completion port.

\begin{table}[htbp]
\centering\small
\caption{Block device command (32 bytes) and completion entry (16 bytes)}\label{tab:blkcmd}
\begin{tabular}{@{}r l l @{\hspace{2.5em}} r l l@{}}
\toprule
\multicolumn{3}{@{}l}{\textbf{Command}} & \multicolumn{3}{l@{}}{\textbf{Completion}}\\\midrule
0 & u8 & opcode: 1 read, 2 write, 3 flush & 0 & u16 & command id\\
1 & u8 & reserved (the buffer returns to its lender) & 2 & u16 & status\\
2 & u16 & command id (echoed) & 4 & u32 & blocks transferred\\
4 & u32 & block count & 8 & u8[8] & reserved, zero\\
8 & u64 & first block (LBA) & & & \\
16 & u64 & byte offset in the buffer capability & & & \\
24 & u8[8] & reserved, zero & & & \\
\bottomrule
\end{tabular}
\end{table}

\begin{table}[htbp]
\centering\small
\caption{Block device status codes}\label{tab:blkstatus}
\begin{tabularx}{\linewidth}{@{}r l X@{}}
\toprule
\textbf{Code} & \textbf{Name} & \textbf{Meaning}\\\midrule
0 & OK & success\\
1 & E\_OP & unknown opcode\\
2 & E\_LBA & more than 256 blocks, or the block range is outside the disk\\
3 & E\_RANGE & offset or length outside the buffer capability, or offset not a multiple of 8\\
4 & E\_PERM & buffer is not T2/T3 memory, or lacks \perm{W} (for read) or \perm{R} (for write)\\
\bottomrule
\end{tabularx}
\end{table}

A command is checked in the order of the codes, and the first failure is reported: the opcode, then the
block range, then the buffer range, then the buffer's space and permission. A command with an unknown
opcode and a bad block range reports E\_OP; a read with a block range beyond the disk into a
capability without \perm{W} reports E\_LBA. A flush is checked for its opcode only. A command that
fails moves no data and reports zero blocks transferred.

\subsection{Protocol}

\begin{figure}[htbp]
\centering
\begin{tikzpicture}[>={Stealth[length=2mm]},font=\small]
  \node[font=\bfseries] (c) at (0,0) {Core};
  \node[font=\bfseries] (d) at (8,0) {Block device (node $N$)};
  \draw[black!40] (0,-0.3) -- (0,-6.2); \draw[black!40] (8,-0.3) -- (8,-6.2);
  \draw[->] (0,-0.8) -- node[above]{\insn{tput} 32-byte command to \hex{2100}} (8,-1.2);
  \draw[->,accent,thick] (0,-1.9) -- node[above]{\insn{cgive c9, N}: buffer capability (= DMA permission)} (8,-2.3);
  \node[anchor=west,align=left,font=\footnotesize] at (8.2,-3.0) {pairs command and capability,\\ moves the data,\\ queues a completion};
  \draw[<-,accent,thick] (0,-4.0) -- node[above]{capability returned to the core's inbox} (8,-3.6);
  \node[anchor=east,font=\footnotesize] at (-0.2,-4.0) {\insn{ctake c9} returns};
  \draw[<-] (0,-5.4) -- node[above]{\insn{tget} 16-byte completion from \hex{2200}} (8,-5.0);
\end{tikzpicture}
\caption{A block-device transfer}\label{fig:blkproto}
\end{figure}

\begin{enumerate}
\item Write the command to the submission port (\insn{tput} of 32 bytes to window offset \hex{100}).
\item Give the buffer capability to the device's node (from the platform table entry), with
  \insn{cgive} or, from user mode, with \insn{csend} through an endpoint for that node.
\item The device pairs a command with a buffer \emph{from the same core that submitted the command},
  and only that core: it records the submitting core of each queued command, and a buffer whose sender
  has no pending command is an \term{orphan} --- a stray gift --- which the device returns to its
  sender unchanged (from the device's node, moving no data) rather than slipping it under another
  core's command. This binds the DMA permission
  to the requester, so one core cannot have its buffer paired with another's command (it matters when
  mutually distrusting software shares the device; see Section~\ref{sec:os-guests}). Among a core's own
  buffers, the oldest is paired first. The device processes a command when it holds a matching buffer
  and has room in its four-entry completion queue; it moves the data between the disk and
  $[\textit{base} + \textit{offset}, \ldots)$ of the capability, appends a completion, and gives the
  capability back to the core that lent it. A \insn{cgive} or \insn{csend} to the device's node
  waits until the core's earlier device writes (the command) have reached the device, so a command
  and the buffer that follows it in program order always meet.
\item \insn{ctake} the capability back, from the device's node only (\insn{ctake cd, rs} with \reg{rs} = $N$). Its arrival means the operation has finished. (A user-mode driver, such as the StrataOS file server, is handed the capability by its supervisor instead.)
\item Read the completion from the completion port and check the command id and status.
\end{enumerate}

The device signals completion the only way a device can, by completing something the core is waiting
for: here, by giving the capability back. A core that has other things to wait for at the same time
waits with \insn{tany} on its inbox.

A block-device read into a capability for framebuffer memory is a write into the framebuffer and
appears on the display like any other.

Every command, including flush, consumes and returns one capability; a flush may be given a
zero-length capability. Reading the completion port when no completion is queued returns 16 zero bytes,
which software can recognise if it never issues command id 0 (the firmware numbers commands from 1). If more than four commands are
submitted without being paired, further commands are discarded. The firmware routines
\texttt{fw\_blk\_read} and \texttt{fw\_blk\_write} implement this protocol (Section~\ref{sec:sfi}).
